\documentclass[10pt,letterpaper]{amsart}
\usepackage[margin=0.84in]{geometry}
\usepackage{tikz}
\usepackage[hidelinks]{hyperref}
\hypersetup{pdftitle={A counterexample to Graffiti.pc Conjecture 430c},pdfauthor={Yuvulmaster}}
\newtheorem*{theorem}{Theorem}
\newtheorem*{lemma}{Lemma}
\title{A counterexample to Graffiti.pc Conjecture 430c}
\author{Yuvulmaster}
\date{}
\begin{document}
\begin{abstract}
We construct line graphs of diameter two whose independent domination numbers are arbitrarily large, while the proposed upper bound is always two. The first example in the family has eleven vertices.
\end{abstract}
\maketitle
An independent dominating set is an independent set such that every vertex
outside it has a neighbor in it. Its minimum size is $i(G)$. Write $M$ for
the maximum-degree vertices of $G$ and $\lambda_{\max}(G)$ for the largest
independence number of a subgraph induced by the neighbors of a single
vertex. The square $G^2$ joins vertices at distance one or two. The
residue $\operatorname{res}$ counts the zeros left at the end of the
Havel--Hakimi process, which repeatedly deletes a largest positive degree
$d$, subtracts one from the next $d$ largest degrees, and reorders.
For connected graphs of order greater than three, Conjecture~430c~\cite{wall}
proposes
\[
 i(G)\le \lambda_{\max}(G)\operatorname{res}(G^2)+\Delta(G[M]).
\]

\begin{theorem}For each $k\ge3$ there is a graph on $k^2+2$ vertices with
independent domination number $k$ and proposed upper bound two.\end{theorem}
\begin{proof}
Start with $K_{k,k}$, choose a vertex $a$ in its left part and a vertex $b$
in its right part, and attach one new leaf to each of $a,b$. Call this
bipartite graph $F$ and let $G$ be its line graph. Thus vertices of $G$
are edges of $F$, and adjacency means sharing an endpoint.

\begin{center}
\begin{tikzpicture}[x=.85cm,y=.53cm,vertex/.style={circle,fill,inner sep=1.5pt}]
\foreach \i/\y in {1/1,2/0,3/-1}{
  \node[vertex] (L\i) at (-1.15,\y) {};
  \node[vertex] (R\i) at (1.15,\y) {};
}
\foreach \i in {1,2,3} {\foreach \j in {1,2,3} {\draw (L\i)--(R\j);}}
\node[above,font=\small] at (L1) {$a$}; \node[above,font=\small] at (R1) {$b$};
\node[vertex] (l) at (-2.15,1) {}; \node[vertex] (r) at (2.15,1) {};
\draw (l)--(L1) (R1)--(r);
\end{tikzpicture}\par
{\footnotesize The graph $F$ for $k=3$. Its eleven edges are the vertices of $G$.}
\end{center}

Independent dominating sets of $G$ correspond exactly to maximal matchings
of $F$. Independence gives a matching, and domination means every unused
edge touches that matching. If some original left vertex and some original
right vertex were both unmatched, their edge could be added. A maximal
matching must therefore saturate all $k$ original vertices on at least one
side, requiring at least $k$ edges. A perfect matching of the original
$K_{k,k}$ is maximal even after the two leaves are added. Thus $i(G)=k$.

Every two disjoint original edges have a cross-edge joining their endpoints.
For a pendant edge and a disjoint original edge, the pendant edge's original
endpoint is adjacent to the opposite-side endpoint of the original edge.
For the two pendant edges, the joining edge is $ab$. Hence every pair of
vertices in $G$ is at distance at most two, and disjoint edges show the
diameter is two. In particular $G^2=K_{k^2+2}$. One Havel--Hakimi step on
the degree sequence of $K_m$ yields the sequence of $K_{m-1}$, so its
residue is one.

The neighbors of a line-graph vertex corresponding to $xy$ are covered by
two cliques, the other edges at $x$ and the other edges at $y$. The two
pendant edges are independent neighbors of $ab$, so $\lambda_{\max}(G)=2$.
Finally, the degree of $xy$ in $G$ is $d_F(x)+d_F(y)-2$. Only $a$ and $b$
have degree $k+1$ in $F$, and $ab$ is their unique joining edge.
Thus $G[M]$ is a single vertex and has maximum degree zero.
The right side is $2\cdot1+0=2$, giving $3\le2$ already at $k=3$.
\end{proof}

\subsection*{Acknowledgment} OpenAI Codex assisted in finding the construction,
checking the 11-vertex example computationally, and preparing the manuscript.
\begin{thebibliography}{1}
\bibitem{wall} E. DeLaVi\~na, \emph{Written on the Wall II: Conjectures of Graffiti.pc},
University of Houston--Downtown, Conjecture~430c (December 8, 2010).
\href{https://www.uhd.edu/academics/sciences/mathematics-statistics/wowii/index.html}{Online catalogue}.
\end{thebibliography}
\end{document}
